\chapter{Conclusion}
GeoCrop addresses the slow, manual and hard-to-scale process of mapping crops by field survey and
GIS digitisation. Research showed that existing tools solve classification, mapping or yield in
isolation, that many studies report optimistic accuracy from random splits, and that research
prototypes become slow when results are shipped as large vector files: our own prototype needed
43.3~MB and 101,435 polygons for one area.

Requirements were defined for a processing pipeline, a plug-in model registry, a spatial catalogue
and a Web GIS portal, with measurable performance targets. The design keeps model results as rasters,
serves only visible tiles through TiTiler and PMTiles, runs heavy work in background workers and
keeps a PostGIS catalogue that drives the interface. Iteration 1 implemented this architecture as
seven containerised services. One screen of AOI1 now needs 337~KB, cold tiles return in 47~ms and
cached tiles in 1~ms, an uploaded Sentinel-2 image is analysed and published in 57~s, and new models
plug in without code changes elsewhere. On the model track, a fair spatial evaluation of eleven
crop-type models showed that a monthly time series is the key input and that a small CPU-only
LightGBM pipeline with neighbourhood inputs and a neighbour vote reaches 76.1\% (AOI1) and 92.2\%
(AOI2), with most remaining errors on field edges.

All platform test cases pass and the main performance targets are met. The main risks are label
noise in the CDL, generalisation to new areas, and limited yield ground truth. FYP-2 continues by
integrating the selected crop-type pipeline, adding field-boundary and yield models, evaluating across
areas, and deploying the hardened system.
